%=============================================================================!
% 10.4  NEIGHBOUR LISTS                                                       !
%=============================================================================!
\section{Neighbour lists}
\label{sec:md-neighbours}

A pair potential cut off at $r_{\mathrm c}$ needs only the pairs closer than
$r_{\mathrm c}$, and in a dense system each atom has a few dozen of them
whatever the size of the system. Testing every pair to find them costs
work that grows as $N^2$, which \cref{sec:pe-cutoffs} promised to reduce
to $N$. This section does so with two lists.

\subsection*{How many pairs matter}

In a liquid of number density $\rho$ (atoms per unit volume), the
neighbours of an atom within $r_{\mathrm c}$ fill a sphere of volume
$\frac43\pi r_{\mathrm c}^3$, so there are about $\frac43\pi r_{\mathrm
c}^3\rho$ of them, and half as many pairs per atom, since each pair has
two atoms. For a dense Lennard-Jones liquid, $\rho = 0.8\sigma^{-3}$, and
$r_{\mathrm c} = 2.5\sigma$, that is 52.4 neighbours and 26 pairs per
atom, against $N(N - 1)/2$ pairs to test. Measured on atoms placed at
random at this density, every $N$ from 500 to 32\,000 gives 26.0 to 26.4
pairs per atom.

\subsection*{The cell list}

The cell is divided into small cells at least $r_{\mathrm c}$ wide: along
the direction of each perpendicular width $w_k$ of \cref{sec:md-image},
there are $n_k = \lfloor w_k/r_{\mathrm c}\rfloor$ of them, each a slice
$1/n_k$ thick in the fractional coordinate $s_k$. Each atom is sorted, in
one pass over the atoms, into the small cell numbered $\lfloor
n_ks_k\rfloor$ along each direction, with $\vb{s}$ wrapped into the cell.
A partner within $r_{\mathrm c}$ changes $s_k$ by less than $r_{\mathrm
c}/w_k \le 1/n_k$ (\cref{eq:md-width}), less than one slice, so it lies
in the atom's own small cell or in one of the $3^3 - 1 = 26$ that touch
it, and only those are searched. A small cell holds on average the same number of
atoms whatever $N$, so the work per atom does not grow with $N$ and the
total grows as $N$. With fewer than three small cells along some
direction, the 26 neighbours would include one cell twice, and
\texttt{mdlab} tests all pairs instead. The method goes back to Quentrec
and Brot\cite{QuentrecBrot1973}.

\Cref{fig:md-neighbours} measures both: testing all pairs grows with a
slope of 2.02 on logarithmic axes, the cell list, from $N = 2000$ on,
with a slope of 0.98. The two cost the same near $N = 1000$, and at $N =
4000$ the cell list is four times faster. At $N = 32\,000$ the cell list
takes \SI{0.78}{\second}; testing all pairs, which takes
\SI{0.42}{\second} at $N = 4000$, would by its slope of 2 take $8^2 = 64$
times as long there, about \SI{27}{\second}.

\begin{figure}[htb]
\centering
\includegraphics{ch10_code/neighbours}
\caption{The time to find every pair within $r_{\mathrm c} = 2.5\sigma$ of
$N$ atoms placed at random at $\rho = 0.8\sigma^{-3}$, by testing all
pairs and through a cell list, against $N$, with lines of slope 2 and 1
(dashed).}
\label{fig:md-neighbours}
\end{figure}

\subsection*{The Verlet list}

Atoms move little in a step, so the pairs within $r_{\mathrm c}$ change
little from one step to the next. Verlet's list\cite{Verlet1967} keeps the
pairs within $r_{\mathrm c} + r_{\mathrm{skin}}$, where the skin
$r_{\mathrm{skin}}$ is a margin added to the cutoff, and reuses them for
many steps, computing at each step only the distances of the listed pairs.
It must be rebuilt before a pair missing from it can come within
$r_{\mathrm c}$. Suppose that since the list was built every atom has
moved less than $\frac12r_{\mathrm{skin}}$, atom $i$ by
$\Delta\vb{r}_i$. A pair not on the list was then at least $r_{\mathrm c}
+ r_{\mathrm{skin}}$ apart. The straight line from $i$'s old position to
$j$'s old position is no longer than the detour through $i$'s new position
and $j$'s new position, since two sides of a triangle are together at
least as long as the third, used once for each triangle of the detour:
\begin{align*}
   r_{ij}^{\text{then}} &\le \lvert\Delta\vb{r}_i\rvert + r_{ij}^{\text{now}}
   + \lvert\Delta\vb{r}_j\rvert ,\\
   r_{ij}^{\text{now}} &\ge r_{ij}^{\text{then}} - \lvert\Delta\vb{r}_i\rvert
   - \lvert\Delta\vb{r}_j\rvert > (r_{\mathrm c} + r_{\mathrm{skin}})
   - \tfrac12r_{\mathrm{skin}} - \tfrac12r_{\mathrm{skin}} = r_{\mathrm c}
   && \text{rearranging.}
\end{align*}
The list is therefore complete until some atom has moved half the skin,
and is rebuilt then. The class \texttt{VerletList} of
\texttt{mdlab.neighbours} keeps the positions at the last build and tests
this at every step:

\begin{code}
    def needs_rebuild(self, positions: ArrayLike) -> bool:
        if self.reference is None:
            return True
        if not np.array_equal(self.cell, self.built_cell):
            return self._strained_needs_rebuild(positions)
        moved = minimum_image(
            np.asarray(positions, dtype=float) - self.reference, self.cell
        )
        return bool(
            np.max(np.einsum("ix,ix->i", moved, moved))
            > (0.5 * self.skin) ** 2
        )
\end{code}

\noindent The second test applies only when the cell itself has changed
since the list was built, as it does under the barostats of
Chapter~14 (\cref{sec:pr-npt}); with the cell fixed the method goes on to
the rule just derived. The displacements are reduced to their nearest
copy so that the
test stays right if the positions have been wrapped back into the cell,
when an atom that crossed a face would otherwise seem to have moved a
whole cell; \texttt{einsum} adds the squares of the three components of
each displacement, and comparing squared lengths with
$(\frac12r_{\mathrm{skin}})^2$ avoids a square root per atom. The list
holds pairs out to $r_{\mathrm c} + r_{\mathrm{skin}}$, so
\cref{eq:md-rule} applies to that distance, $r_{\mathrm c} +
r_{\mathrm{skin}} \le \frac12w_{\min}$, and \texttt{mdlab} refuses a
smaller cell. A thicker skin means fewer rebuilds and more pairs to check
at each step. In the argon run of \cref{sec:md-loop}, with a skin of
\SI{1}{\angstrom}, the list is built 187 times in 2000 steps, about every
eleventh step, each time by testing all pairs: its box of
\SI{26.3}{\angstrom} holds only $\lfloor26.3/9.5\rfloor = 2$ small cells
of $r_{\mathrm c} + r_{\mathrm{skin}} = \SI{9.5}{\angstrom}$ along each
side.

\begin{keyidea}
Each atom has about $\frac43\pi r_{\mathrm c}^3\rho$ neighbours within the
cutoff. A cell list finds them by searching only neighbouring small cells,
at a cost that grows as $N$; a Verlet list with a skin
$r_{\mathrm{skin}}$ reuses the pairs within $r_{\mathrm c} +
r_{\mathrm{skin}}$ until some atom has moved $\frac12r_{\mathrm{skin}}$.
\end{keyidea}

\notebook{10}{time all pairs against the cell list}

\begin{selfcheck}
With $r_{\mathrm c} = 2.5\sigma$ and a skin of $0.3\sigma$, how far may
the atoms move before the Verlet list must be rebuilt, and about how many
pairs per atom does the list hold at $\rho = 0.8\sigma^{-3}$?
\end{selfcheck}
